The topological description is in terms of basic
hardware units such as registers, selectors,
control structures etc. linked by signals or signal
arrays.

The topological description class includes the following
fields -

{\bf type}

This indicates the type of element. Types are (see class
{\bf TDEConstants} -
\blist
\item   TDEType.SIG -       declare a signal or signal array
\item   TDEType.CONNECT -   connect a signal or signal array to another or to a
                            constant
\item   TDEType.SELECT -    select one of several inputs
\item   TDEType.TSELECT -   a 3-state driver
\item   TDEType.REG -       a register
\item   TDEType.WAIT -      wait for completion of parallel execution streams
\item   TDEType.DIVERGE -   handle queue divergence
\item   TDEType.INV -       invert a signal
\item   TDEType.AND -       AND 2 or more signals
\item   TDEType.OR -        OR 2 or more signals
\item   TDEType.XOR -       XOR 2 or more signals
\item   TDEType.OPERATOR -  apply an expression operator to operands
\item   TDEType.START -     start an execution stream
\item   TDEType.EXECP -     delay execution on queue availability
\item   TDEType.ILOOP -     control an infinite loop
\item   TDEType.DEL -       delay a signal one or more clock cycles
\item   TDEType.WHEN -      a 'when' statement
\item   TDEType.WHILE -     a 'while' statement
\item   TDEType.DOWHILE -   a 'do while' statement
\item   TDEType.PRIORITY -  a priority encoder
\item   TDEType.RESYNC -    resynchronise a pulse or edge to another clock
\item   TDEType.SAMPLE -    sample a value in another clock domain
\item   TDEType.QUEUEBUFFER - a queue variable (buffer)
\item   TDEType.IBUF -      an FPGA pad input buffer
\item   TDEType.OBUF -      an FPGA pad output buffer
\item   TDEType.DFF -       a D-type flip-flop
\item   TDEType.CRAM -      a combinatorial-output RAM
\item   TDEType.RRAM -      a registered-output RAM
\item   TDEType.CLOCK -     create a clock signal source
\item   TDEType.CLOCKINFO - a wrapper for information about a clock signal
\item   TDEType.ELEMENT -   create a specific FPGA vendor library element
\item   TDEType.PORT -      an external signal connection to/from another netlist
\item   TDEType.SPECIAL -   special functions for specific FPGA families
%            \item   TDEType.SIMREAD -   read a value for a variable from a file during simulation
%            \item   TDEType.SIMWRITE -  write a value from a variable to a file during simulation
%            \item   TDEType.SIMVAR -    communicate program variables to the simulator
\blend

{\bf ArrayList params} is a parameter section which includes such
things as initial values and flags. Items in the parameter
list may be type {\bf Integer}, {\bf String} or {\bf
Boolean}.

{\bf ArrayList inputs} and {\bf ArrayList outputs} are lists of 
TDEVars which are inputs and outputs to the TDE.

{\bf location}

This field gives the source file location.
Since these structures represent the hardware rather than the
source code the source file location does not provide a very
good cross reference to the source code. For example the
structure for a variable declaration includes all the logic
required to assign values to that variable, including input
selectors, however the selector logic is derived from the
various assignment statements which target that variable. The
line numbers of these statements are not available in the
variable structure. We could partition the logic into smaller
units to improve this situation, but that might impose
restrictions on the implementation and might also preclude some
simple low-level optimisations.

{\bf active}

This field indicates if the TDE is still in use.
During optimisation some TDEs are eliminated
and rather than delete them from the list of elements they
are instead tagged as inactive via this field.

{\bf no\_delete}

If this field is true the element should not be deleted. I
think the only use for this is for the SELECT TDE where
it is required when used for queue variable assignment that
a single input SELECT not be eliminated as redundant.

\subsection{constructors}

    There are a number of constructors all of which require the
    TDE type but which may omit parameter, input or output lists
    or source file location. Where omitted the lists are added
    later.

\subsection{methods}


    {\bf add2i()} \\ Add a signal to the input list.

    {\bf add2o()} \\ Add a signal to the output list.

    {\bf add2p()} \\ These add an item to the parameter list. There are multiple
    signatures for different argument types.

    {\bf addPin2Element()} \\ This method adds a pin and a connected signal to a
    TDE instance of type ELEMENT.

    {\bf addProperty2Element()} \\ This method adds a property to a TDE instance
    of type ELEMENT. The property name and value are added to the parameter list
    of the TDE instance.

    {\bf setElementPin()} This method adds a pin to a TDE instance of type
    ELEMENT. The port name, type (input, output, 3-state output or clock),
    associated TDEVar, port width and EDIF port array format are passed as
    arguments. The array size is used later (in XTDECode) to pad input nets
    formed from the TDEvar so that the full port width is connected even though
    the signal array may be narrower. Padding is GND unless the TDEVar is for a
    signed type in which case the MSB is propagated.


    {\bf redundant()} \\ This method examines a TDE to see if it contains
    redundancies or is entirely redundant. For example a duplicated selector
    data input is removed and the select input is ORed with the duplicate
    select. Logic gates with a single input are effectively removed (see below)
    and the single input is connected to the output directly.

    If the TDE is completely redundant it has all inputs, outputs and parameters
    removed, the active flag is made false and the method returns true. If the
    TDE is already in the TDE list it will be recognised as inactive during any
    traversals of the list and will be ignored.

    The method is called in only two places, TDEList.addTDE() and
    TDEList.optimise(). The first is called to add a TDE to the TDE list but
    does not do so if the TDE is wholly redundant (redeundant() returns true).
    The second is a TDE list optimiser which calls redundant() only for SELECT
    TDEs as a part of its optimisation of the list. Since the TDE is already in
    the list a true return from redundant() simply means the TDE has been made
    inactive and will be ignored when encountered again in any context.


    {\bf ncode()} \\ This method generates netlist code from the TDE. A switch
    on the TDE type selects a matching method in abstract class {\bf TDECode},
    these methods being overridden in class XTDECode in the case of Xilinx.

    {\bf deactivate()} \\ Make a TDE inactive by unlinking all inputs and
    outputs from the associated TDEVars (removing the reverse link), clearing
    the parameter list and setting boolean {\bf active} to false.

    {\bf removeInput()} \\ Remove a control signal input to a TDE. This is not
    used for signals which are arrays. The TDEVar is removed from the list of
    inputs and this TDE is removed from the destination list of the associated
    TYDEVar. This method assumes that there is only one occurrence of the TDEVar
    in the input list.

    {\bf replaceInput()} \\ Replace a TDEVar in the input list with another. The
    source link in the associated TDEVar is not changed by this method; it is
    assumed that that is done elsewhere.

    {\bf replaceOutput()} \\ Replace a TDEVar in the output list with another.
    The destination link in the associated TDEVar is not changed by this method;
    it is assumed that that is done elsewhere.
    
    {\bf dump()} \\ print a description of the TDE to the .icl file.


\subsection{signal linking}

    Input and output signals to a TDE are representad by class TDEVar.
    TDE fields {\bf input} and {\bf output} are lists of TDEVars (class
    {\bf TDEVarList} which extends class {\bf ArrayList}). The function of
    each input and output is dependant on the specific TDE type. A TDEVar
    in turn has a list of TDEs to which it is linked (within subclass Common).
    
    Optimisation of TDEs must handle this double linking where TDEVars are
    to be removed or relocated.
    
